\documentclass{exam}

\newif\ifA
\newif\ifkey
\newif\ifprintselected   % required by exam_config's \ansbox

%%%%%%%%%%%%%%%%%%%%%%%%%%%%%%%%%%%%%%%%%%%%%%%%%%%%%%
%%% SET THESE IF VARIABLES
%%%%%%%%%%%%%%%%%%%%%%%%%%%%%%%%%%%%%%%%%%%%%%%%%%%%%%
% Lab name -- labs are named by the lesson they belong to, never by a
% sequential lab number (Neil, 2026-10-12): this is the L12 Lab.
\def\labLesson{12}

% is this the key or not?
\ifdefined\iskey
	\ifnum\iskey=1
		\keytrue
	\else
		\keyfalse
	\fi
\else
	%% Pick one manually
	\keyfalse % if this is NOT the key
%	\keytrue     % if this IS the key
\fi

% set up the header and footer - change for every term
\ifdefined\thisTermIn
	\def\thisterm{\thisTermIn}
\else
	\def\thisterm{Spring 2027}
\fi
%
\def\thisclass{ECE 448}

%%%%%%%%%%%%%%%%%%%%%%%%%%%%%%%%%%%%%%%%%%%%%%%%%%%%%%
%%% END SET THESE IF VARIABLES
%%%%%%%%%%%%%%%%%%%%%%%%%%%%%%%%%%%%%%%%%%%%%%%%%%%%%%
% packages and shortcuts
\input{myPackages_gr}
\usepackage[hidelinks, linkcolor=red]{hyperref}
\input{myShortcuts}
\setlength{\parindent}{0pt}
\setlength{\parskip}{6pt}

%% exam class customization
\input{exam_config}
% ECE 448 document fonts.
%
% Body text is Barlow (SIL OFL; the faces and the license are in
% latex/fonts/Barlow/). Barlow is not in TeX Live, so it is loaded as an
% OpenType face through fontspec -- which means the document must be built
% with lualatex, not pdflatex. build_practice.sh and build_lab.sh both do.
%
% Compiling with pdflatex anyway is not an error: the guard below falls back
% to the LaTeX default so the source still compiles under any engine and can
% be syntax-checked. The resulting PDF just is not the one we ship.
%
% Paths are relative to latex/, which is where both build scripts run
% (each does cd "$(dirname "$0")"). Loading by path rather than by installed
% family name means no system font install and no fontconfig, so a fresh
% clone builds a byte-comparable PDF anywhere.
%
% MATH IS DELIBERATELY LEFT IN COMPUTER MODERN. Barlow has no math companion,
% and CM math under sans body text is already the house pairing: the decks set
% words in Source Sans Pro and let MathJax render symbols in its CM-derived
% face. Moving math to a sans as well would mean unicode-math plus a sans math
% font (Fira Math being the plausible one) and re-checking every equation in
% the course for symbol coverage. That is a separate decision, not a side
% effect of this one.

\usepackage{iftex}
\ifLuaTeX
	\usepackage{fontspec}
	\defaultfontfeatures{Ligatures=TeX}
	\setmainfont{Barlow}[
		Path           = fonts/Barlow/,
		Extension      = .ttf,
		UprightFont    = *-Regular,
		ItalicFont     = *-Italic,
		BoldFont       = *-Bold,
		BoldItalicFont = *-BoldItalic,
	]
	\setsansfont{Barlow}[
		Path           = fonts/Barlow/,
		Extension      = .ttf,
		UprightFont    = *-Regular,
		ItalicFont     = *-Italic,
		BoldFont       = *-SemiBold,
		BoldItalicFont = *-BoldItalic,
	]
	% microtype earns its place here specifically because of Barlow. Its
	% glyphs are narrower than Computer Modern's, so justified prose sets
	% tighter and TeX has to squeeze interword space to fit; two paragraphs
	% in the L02/L03 sets went overfull on the first Barlow build. Protrusion
	% and font expansion absorb that -- both are lualatex features, so this
	% belongs inside the guard with everything else.
	\usepackage{microtype}

	% Mono stays Latin Modern: Barlow ships no monospaced face, and verbatim
	% and code need one that actually is fixed-pitch.
\fi
   % Barlow body text (needs lualatex)
\printselectedfalse      % full worked solutions in the key, not answers-only
\CorrectChoiceEmphasis{\color{red}}
\SolutionEmphasis{\color{red}}
\renewcommand\questionlabel{\textbf{\thequestion.}}
\renewcommand{\thepartno}{\alph{partno}}
\renewcommand{\partlabel}{(\thepartno)}

% Fillable table cell: blank in the student copy, the expected value in red on
% the KEY. Fixed-width centered columns keep the blank cells writable.
\newcolumntype{K}[1]{>{\centering\arraybackslash}p{#1}}
\newcommand{\kv}[1]{\ifkey{\color{red}#1}\fi\rule[-0.10in]{0pt}{0.34in}}

% print the answers if this is the key
\ifkey
	\printanswers % if this IS the key
\else
	\noprintanswers  %if this is NOT the key
\fi

% print the header and footer
% need to print key on top if this is the key
\ifkey
	\firstpageheader{}{\color{red}{ KEY}}{}
	\runningheader{\thisclass\hspace{1pt} -- \thisterm }{L\labLesson\ Lab\hspace{1pt}\color{red}{ KEY}}{\thepage\hspace{1pt} of\hspace{1pt} \numpages}
	\firstpagefooter{}{}{}
	\runningfooter{}{}{}
\else
	\firstpageheader{}{}{}
	\runningheader{\thisclass \hspace{1pt} -- \thisterm }{L\labLesson\ Lab\hspace{1pt}}{\thepage\hspace{1pt} of\hspace{1pt} \numpages}
	\firstpagefooter{}{}{}
	\runningfooter{}{}{}
\fi

\runningheadrule


\begin{document}
\pagestyle{headandfoot}   % exam class

%%%%%%%%%%%%%%%%%%%%%%%%%%%%%%%%%%%%%%%%%%%%%%%%%%%%%%
% title block
%%%%%%%%%%%%%%%%%%%%%%%%%%%%%%%%%%%%%%%%%%%%%%%%%%%%%%

\begin{center}
USAF ACADEMY \\
DEPARTMENT OF ELECTRICAL AND COMPUTER ENGINEERING \\
LESSON \labLesson\ LAB - CW \\
\textbf{Lab Packet} \\
\thisterm \\[4mm]
\end{center}

\vspace{-4pt}
\textbf{Name:} \rule[-2pt]{2.6in}{0.4pt} \hfill
\textbf{Section:} \rule[-2pt]{0.9in}{0.4pt} \hfill
\textbf{Date:} \rule[-2pt]{1.1in}{0.4pt}

\vspace{4pt}

\fullwidth{\textbf{Learning Objective 6.1: } I can transmit and receive CW (Morse) using
GNU Radio and decode a received CW stream.}

\vspace{2pt}

This packet is the L\labLesson\ Lab, in three parts, and you turn it in once, complete,
before Lesson 13. Do \textbf{Part 1} on your own before the lab period: run the link
flowgraph and make the predictions you will check at the bench. \textbf{Part 2} is
your own measurements, recorded in class: Questions 3 to 5 on your own laptop, and
Question 6 at a PlutoSDR station with your team. Do \textbf{Part 3} on your own after
the lab: compare the measurements with the predictions and answer the analysis
questions.

\textbf{Grading.} You submit the packet once, we grade it against Objective 6.1 on
the course's 1 / 0 scale, and you may revise it under course policy. Paste screenshots where a
question asks for one.

\textbf{Transmitting.} In class we transmit only into a cable, into a shielded
enclosure, or at very low power in an ISM band, and only when the lesson says to.
Questions 1 to 5 run on your laptop and transmit nothing. Question 6 transmits from a
PlutoSDR's TX port, through a cable with an inline 30 dB SMA attenuator, to the same
Pluto's RX port, at 915 MHz. Connect the cable and attenuator before you run the
flowgraph, and never connect an antenna in this lab. Never set TX attenuation below
30 dB or change the slider's range; the RX input is rated to only a few dBm. We have
run the PlutoSDR flowgraph only in simulation, so Question 6 is its first run on
hardware.

The flowgraphs are on the L12 lesson page, \texttt{L12-cw-tx.grc} and
\texttt{L12-cw-link.grc}, and on the L12 Lab page, \texttt{L12-cw-pluto.grc}. Every
number below assumes their settings: the keyed carrier at 5 kHz in a 32 kHz band,
noise spread evenly over that band, a 700 Hz BFO, 5 ms edges, and the default message
\texttt{USAFA ECE 448}. The SNR in a CW filter of width $B$ at noise amplitude $\sigma$
is
\[ \text{SNR} = \frac{32{,}000}{\sigma^2 B}. \]

\begin{questions}

\fullwidth{\large\textbf{Part 1: Pre-Lab} \normalsize(individual, before the lab period)}

%%%%%%%%%%%%%%%%%%%%%%%%%%%%%%%%%%%%%%%%%%%%%%%%%%%
\question \textbf{Run the link.} Open \texttt{L12-cw-link.grc} in GRC and run it
with every slider at its default.

\begin{parts}

	\begin{minipage}[t]{0.58\linewidth}
		\part After about 30 s, record the Received label, and time one repeat of the
		message with a watch.
		\begin{solution}[0.7in]
		The label reads \texttt{USAFA ECE 448} three times with no errors. One repeat, 126
		units at 60 ms, takes 7.56 s.
		\end{solution}
	\end{minipage}\hfill%
	\begin{minipage}[t]{0.37\linewidth}
		\raggedleft \vspace{12pt}
		\ansbox[1.4in]{no errors, 7.6 s}
	\end{minipage}

	\part Paste a screenshot of the running flowgraph's window.
		\begin{solution}[1.6in]
		The waterfall shows the keyed carrier at +5 kHz, the envelope sits near 0.64 during
		marks and near 0.04 between them, and the key trace follows the marks.
		\end{solution}

\end{parts}

%%%%%%%%%%%%%%%%%%%%%%%%%%%%%%%%%%%%%%%%%%%%%%%%%%%
\newpage
\question \textbf{Predictions.} You will check each of these at the bench.

\begin{parts}

	\begin{minipage}[t]{0.58\linewidth}
		\part Find the length of a dit, a dah, the gap between letters, and the gap
		between words at 13 WPM.
		\begin{solution}[1.1in]
		\eq{ t_\text{dit} &= 1.2\ \text{s} / 13 = 92.3\ \text{ms} \\
			t_\text{dah} &= t_\text{letter} = 3 \times 92.3 = 277\ \text{ms} \\
			t_\text{word} &= 7 \times 92.3 = 646\ \text{ms} }
		\end{solution}
	\end{minipage}\hfill%
	\begin{minipage}[t]{0.37\linewidth}
		\raggedleft \vspace{12pt}
		dit = \ansbox[1.0in]{92.3 ms}

		\vspace{4pt}
		dah = \ansbox[1.0in]{277 ms}

		\vspace{4pt}
		word gap = \ansbox[1.0in]{646 ms}
	\end{minipage}

	\part Write your initials, two or three letters, in Morse code. Then write them as
	a list with one 0 or 1 per unit: 1 for each unit the key is down, 0 for each unit
	it is up, ending with the 7-unit gap after the word. Count the units and find how
	long the pattern lasts at 13 WPM.
		\begin{solution}[1.2in]
		The answer depends on the initials. For JD, J is \texttt{.---} and D is
		\texttt{-..}, so the list is
		\begin{center}\texttt{[1,0,1,1,1,0,1,1,1,0,1,1,1, 0,0,0, 1,1,1,0,1,0,1, 0,0,0,0,0,0,0]}\end{center}
		J takes 13 units, the letter gap 3, D 7, and the word gap 7, for 30 units in all,
		or $30 \times 92.3\ \text{ms} = 2.77$ s.
		\end{solution}

	\part Predict the SNR in the CW filter at each setting you will measure in
	Question 4.

	\begin{center}
	\renewcommand{\arraystretch}{1.25}
	\begin{tabular}{| l | K{0.75in} | K{0.75in} | K{0.75in} | K{0.75in} | K{0.75in} |}
	\hline
	$\sigma$ & $4.5$ & $5.5$ & $6.5$ & $6.5$ & $6.5$ \\\hline
	Filter width & 200 Hz & 200 Hz & 200 Hz & 1000 Hz & 50 Hz \\\hline
	SNR (dB) & \kv{$9.0$} & \kv{$7.2$} & \kv{$5.8$} & \kv{$-1.2$} & \kv{$11.8$} \\\hline
	\end{tabular}
	\end{center}

		\begin{solution}[0.4in]
		For example, $32{,}000/(4.5^2 \times 200) = 7.90$, which is 9.0 dB.
		\end{solution}

	\part The decoder is set to 20 WPM, so its unit is 60 ms. It calls a mark under
	2 units a dit, ends a letter when a space reaches 2 units, and ends a word when a
	space reaches 5. A sender at $W$ WPM keys perfect timing. Find the slowest and the
	fastest $W$ the decoder should copy, and say which window sets each limit.
		\begin{solution}[1.6in]
		The sender's dit is $d = 1.2/W$ s. Its letter gap, $3d$, must be at least 120 ms
		and under 300 ms, and its word gap, $7d$, must reach 300 ms. Its dit must stay
		under 120 ms and its dah, $3d$, must reach 120 ms.
		\eq{ 3d < 300\ \text{ms} &\Rightarrow d < 100\ \text{ms},\ W > 12 \\
			7d \ge 300\ \text{ms} &\Rightarrow d \ge 42.9\ \text{ms},\ W \le 28 }
		The other windows allow a wider range. Below 12 WPM every letter gap reads as a
		word gap; above 28 WPM the word gaps read as letter gaps. Both limits sit on a
		boundary, so the decoder may fail at exactly 12 and 28 WPM, where a gap falls on
		the 5-unit boundary.
		\end{solution}

	\begin{minipage}[t]{0.58\linewidth}
		\part The CW filter is 200 Hz wide and centered on the 700 Hz tone. How far can
		you tune the receiver from the carrier before the tone leaves the filter?
		\begin{solution}[0.6in]
		You can tune half the filter width, 100 Hz, either way: the tone moves from 700 Hz to
		$700 \pm 100$ Hz, the edges of the filter.
		\end{solution}
	\end{minipage}\hfill%
	\begin{minipage}[t]{0.37\linewidth}
		\raggedleft \vspace{12pt}
		\ansbox[1.1in]{100 Hz}
	\end{minipage}

\end{parts}

%%%%%%%%%%%%%%%%%%%%%%%%%%%%%%%%%%%%%%%%%%%%%%%%%%%
\newpage
\fullwidth{\large\textbf{Part 2: Measurements} \normalsize(in class)}

\question \textbf{Build a hand keyer.} Start GRC with an empty canvas. Do not open
\texttt{L12-cw-tx.grc}: build the keyer yourself as the lab page lists, with your
list from Question 2(b) in the Vector Source.

\begin{parts}

	\part Paste a screenshot of your flowgraph on the GRC canvas. Every connection
	must be black, not red.
		\begin{solution}[1.5in]
		The flowgraph matches the lab page's list: Vector Source, Repeat, and Throttle, then a
		Moving Average, with a two-input QT GUI Time Sink and a two-input QT GUI
		Frequency Sink fed by the Throttle and the Moving Average.
		\end{solution}

	\begin{minipage}[t]{0.58\linewidth}
		\part Paste a screenshot of the time sink zoomed in on one dit, and compare its
		length with Question 2(a).
		\begin{solution}[1.3in]
		The Repeat block holds each unit for $\texttt{int}(8000 \times 1.2/13) = 738$
		samples, 92.25 ms, within a fraction of a millisecond of 92.3 ms. The shaped
		trace rises in a straight 5 ms ramp and lags the hard trace by 2.5 ms, half the
		Moving Average's length.
		\end{solution}
	\end{minipage}\hfill%
	\begin{minipage}[t]{0.37\linewidth}
		\raggedleft \vspace{12pt}
		\ansbox[1.1in]{92 ms}
	\end{minipage}

	\part Read the highest peak of the hard and shaped spectra within 100 Hz of 1 kHz and of
	2 kHz from 0 Hz, where the carrier would sit. How far below hard keying is your ramp at each? Lesson 12's 5 ms raised
	cosine was 40 dB and 53 dB below; explain the difference.

	\begin{center}
	\renewcommand{\arraystretch}{1.25}
	\begin{tabular}{| l | K{1.1in} | K{1.1in} | K{1.1in} |}
	\hline
	& Hard (dB) & Ramp (dB) & Difference (dB) \\\hline
	1 kHz from the carrier & & & \kv{$\approx -25$} \\\hline
	2 kHz from the carrier & & & \kv{$\approx -30$} \\\hline
	\end{tabular}
	\end{center}

		\begin{solution}[1.0in]
		The ramp is about 23 to 26 dB below hard keying 1 kHz away and 29 to 31 dB below
		2 kHz away. A straight ramp still has a sharp corner where it starts and where
		it ends, and those corners spread energy farther than the smooth start and end
		of a raised cosine of the same length, which is 15 to 23 dB lower.
		\end{solution}

\end{parts}

%%%%%%%%%%%%%%%%%%%%%%%%%%%%%%%%%%%%%%%%%%%%%%%%%%%
\newpage
\question \textbf{The receiver in noise.} Follow the lab page's steps for
\texttt{L12-cw-link.grc}. At each setting, set the threshold halfway between the space
and mark levels you read on the envelope plot, wait for three new repeats of the
message, and count the wrong, missing, and extra characters among those 39.

\begin{parts}

	\part Fill in the table.

	\begin{center}
	\renewcommand{\arraystretch}{1.25}
	\begin{tabular}{| l | K{0.75in} | K{0.75in} | K{0.75in} | K{0.75in} | K{0.75in} |}
	\hline
	$\sigma$ & $4.5$ & $5.5$ & $6.5$ & $6.5$ & $6.5$ \\\hline
	Filter width & 200 Hz & 200 Hz & 200 Hz & 1000 Hz & 50 Hz \\\hline
	Space level & \kv{$0.19$} & \kv{$0.24$} & \kv{$0.28$} & \kv{$0.62$} & \kv{$0.14$} \\\hline
	Mark level & \kv{$0.67$} & \kv{$0.68$} & \kv{$0.69$} & \kv{$0.86$} & \kv{$0.64$} \\\hline
	Threshold & \kv{$0.43$} & \kv{$0.46$} & \kv{$0.48$} & \kv{$0.74$} & \kv{$0.39$} \\\hline
	Wrong of 39 & \kv{0} & \kv{1 to 3} & \kv{$\approx 5$} & \kv{$\approx 21$} & \kv{0 or 1} \\\hline
	\end{tabular}
	\end{center}

		\begin{solution}[0.8in]
		Averaged over four runs of 600 characters, the wrong characters were 0.2\%, 4.3\%,
		14\%, 55\%, and 0.8\%. Counts over 39 characters vary from run to run because the
		noise is random. The space level rises with the noise and with the filter width,
		so the threshold has to move with it.
		\end{solution}

	\part Paste a screenshot at $\sigma = 6.5$ with the 200 Hz filter, and mark one wrong
	character in the Received label.
		\begin{solution}[1.4in]
		Answers vary. Each error comes from a mark that noise splits, shortens, or invents,
		or from a lost dit: a dah broken into two dits turns A into S; a dah shrunk to a dit
		turns 8 into 7; a lost dit turns F into R or into an unknown pattern \#; and a noise
		burst adds an extra E.
		\end{solution}

\end{parts}

%%%%%%%%%%%%%%%%%%%%%%%%%%%%%%%%%%%%%%%%%%%%%%%%%%%
\newpage
\question \textbf{Speed and tuning.} Follow the lab page's steps, with the noise
amplitude at 1, the 200 Hz filter, and the threshold at 0.3.

\begin{parts}

	\part With the decoder at 20 WPM, record whether each sender speed decodes
	correctly, and if not, what goes wrong.

	\begin{center}
	\renewcommand{\arraystretch}{1.25}
	\begin{tabular}{| l | K{0.5in} | K{0.5in} | K{0.5in} | K{0.5in} | K{0.5in} | K{0.5in} | K{0.5in} |}
	\hline
	Sender WPM & 10 & 12 & 13 & 20 & 27 & 28 & 30 \\\hline
	Decodes? & \kv{no} & \kv{no} & \kv{yes} & \kv{yes} & \kv{yes} & \kv{no} & \kv{no} \\\hline
	\end{tabular}
	\end{center}

		\begin{solution}[1.2in]
		At 10 WPM every dit lasts 2 decoder units, so the decoder reads every mark as a dah
		and prints Ts and Ms. At 12 WPM each letter gap lasts exactly 5 units and reads as a
		word gap, so most letters come out separated by spaces, as in \texttt{US A F A}. From 13 to
		27 WPM the text is correct. At 28 WPM each word gap lasts exactly 5 units, and the
		words run together; at 30 WPM the letter gaps also fall to the 2-unit boundary, the
		letters merge, and the decoder prints \#.
		\end{solution}

	\part Record the pitch and whether the text decodes at each tuning.

	\begin{center}
	\renewcommand{\arraystretch}{1.25}
	\begin{tabular}{| l | K{0.9in} | K{0.9in} | K{0.9in} | K{0.9in} |}
	\hline
	Tuning (Hz) & 5050 & 5080 & 5100 & 5150 \\\hline
	Pitch & \kv{lower, 650 Hz} & \kv{lower, 620 Hz} & \kv{lower, 600 Hz} & \kv{lower, 550 Hz} \\\hline
	Decodes? & \kv{yes} & \kv{yes} & \kv{no} & \kv{no} \\\hline
	\end{tabular}
	\end{center}

		\begin{solution}[0.9in]
		Tuning above the carrier puts it below 0 Hz after the Xlating filter, so the tone
		drops by the tuning error. At 100 Hz the tone sits on the filter's edge, the mark
		level falls to about half, and the text breaks up; at 150 Hz the filter removes
		the tone and nothing decodes.
		\end{solution}

\end{parts}

%%%%%%%%%%%%%%%%%%%%%%%%%%%%%%%%%%%%%%%%%%%%%%%%%%%
\newpage
\question \textbf{CW on a PlutoSDR.} At a PlutoSDR station, with your team, follow the
lab page's steps: connect the Pluto's TX port to its RX port through the 30 dB
attenuator, check the Pluto with \texttt{iio\_info -s}, and run
\texttt{L12-cw-pluto.grc}. The transmitter and receiver share the Pluto's clock, so
the carrier should arrive close to +5 kHz. When you finish, stop the flowgraph and
unplug the Pluto before you disconnect the cable.

\begin{parts}

	\part Once \texttt{USAFA ECE 448} decodes, paste a screenshot that shows the
	waterfall, the envelope with the threshold line halfway between the space and mark
	levels, and the Received label.
		\begin{solution}[1.5in]
		The waterfall shows the keyed carrier at +5 kHz. The AGC holds the envelope's mark level
		near 0.6, the red threshold line sits halfway between the space and mark levels, and
		the message decodes without errors.
		\end{solution}

	\part Record your teammate's message as decoded, and check it with them.
		\begin{solution}[0.6in]
		The decoded text matches the teammate's message. The keyer sends lowercase letters
		as capitals and skips any character that is not in its Morse table, such as @.
		\end{solution}

	\part Record the mark and space levels and whether the text decodes at each TX
	attenuation. Explain why the mark level hardly changes while the received power falls
	by 10 dB a step, and what does change.

	\begin{center}
	\renewcommand{\arraystretch}{1.25}
	\begin{tabular}{| l | K{0.9in} | K{0.9in} | K{0.9in} | K{0.9in} |}
	\hline
	TX attenuation (dB) & 60 & 70 & 80 & 89 \\\hline
	Mark level & & & & \\\hline
	Space level & & & & \\\hline
	Decodes? & & & & \\\hline
	\end{tabular}
	\end{center}

		\begin{solution}[1.0in]
		Each 10 dB step divides the received amplitude by $10^{10/20} = 3.16$, and the AGC
		raises its gain by the same factor, so the mark level stays near 0.6. The noise
		does not fall with the signal, so the AGC raises it too, and the space level climbs
		as the SNR drops. With estimated Pluto levels, a simulation decoded at every step,
		with the space level rising from about 0.04 to 0.05; the real levels, and whether
		the text still decodes at 89 dB, depend on the hardware.
		\end{solution}

\end{parts}

%%%%%%%%%%%%%%%%%%%%%%%%%%%%%%%%%%%%%%%%%%%%%%%%%%%
\newpage
\fullwidth{\large\textbf{Part 3: Writeup} \normalsize(individual, after the lab)}

\question \textbf{Analysis.}

\begin{parts}

	\part Compare your noise table with Question 2(c). Find the lowest SNR in the 200 Hz
	filter at which the link copied without errors, explain why the threshold had to
	move as the noise rose, and say what SNR a Lesson 2 link budget for this link
	should deliver in the CW filter.
		\begin{solution}[1.5in]
		The link copied without errors at 9.0 dB ($\sigma = 4.5$), lost about 1 in 25 at
		7.2 dB, and lost about 1 in 7 at 5.8 dB. Noise passing the filter raises
		the envelope during the spaces, from about 0.04 at $\sigma = 1$ to 0.28 at
		$\sigma = 6.5$, so a fixed threshold of 0.3 lets noise peaks through as dits; the
		threshold must sit halfway between the space and mark levels. A budget should
		deliver at least 10 dB in the 200 Hz filter, plus a fade margin.
		\end{solution}

	\part Compare the 1000, 200, and 50 Hz filters at $\sigma = 6.5$. Explain why the
	narrowest filter did best, and use your tuning measurement to say what it gives
	up.
		\begin{solution}[1.5in]
		The signal is only about 100 Hz wide, so a wider filter adds noise without adding
		signal: the 1000 Hz filter passed five times the noise of the 200 Hz filter and lost
		more than half the characters, while the 50 Hz filter passed a quarter of it and lost
		under 1\%. The tone must stay
		inside the filter: the 200 Hz filter decoded out to 80 Hz of tuning error and
		failed at 100 Hz, so a 50 Hz filter would need the receiver within about 25 Hz of
		the carrier, and transmitter drift would matter four times as much.
		\end{solution}

	\part Compare the sender speeds that decoded with Question 2(d), say which window
	failed at each end, and propose a change to the decoder that would let it copy any
	sender's speed.
		\begin{solution}[1.5in]
		The prediction was $12 < W \le 28$, and the link decoded from 13 to 27 WPM.
		At the slow end the letter-gap window failed first, as letter gaps reached the
		5-unit word boundary; at the fast end the word gaps fell to that same boundary.
		At both edge speeds a gap lands exactly on the 5-unit boundary, and the noise's
		jitter of a few hundredths of a unit pushes it to either side, so the text breaks up. A decoder can measure the sender's speed instead of
		assuming it: keep a running average of the short marks, take it as the dit
		length, and set the 2- and 5-unit boundaries from that.
		\end{solution}

\end{parts}

\end{questions}

\vspace{1.0em}
\noindent\textbf{Documentation:} \rule[-2pt]{0.6\linewidth}{0.4pt}

\end{document}
